\documentclass[main.tex]{subfiles}


\begin{document}

% \AddToShipoutPictureFG{%
% \begin{tikzpicture}[overlay,remember picture]
% \path (current page.south west) -- (current page.north east)
% node[midway,scale=1,color=gray,sloped,opacity=0.4] {\textnormal{Кравченко Игорь Игоревич\hspace{1cm} +79010144910\hspace{1cm} super.antig@yandex.ru}};
% \end{tikzpicture}
% }

% %from pagenumber to up
% \phantomsection 
% \hypertarget{toc}{}

 

 

% номер секции вручную
\setcounter{section}{8
}
\addtocounter{section}{-1} 

\section{Средняя скорость}


При неравномерном движении скорость\footnote{И ускорение тоже.} в общем случае может меняться. Это имеет место, например, при непредсказуемом движении муравья (рис.\,\ref{fig:p16}).

\begin{figure}[H]
    \centering

    \begin{tikzpicture}[font=\small,            line cap=round,                             >={Stealth[inset=0pt,round,length=3mm, angle'=20]},vec/.style={>={Stealth[inset=0pt,round,length=3.5mm, angle'=20]}}]


    \draw [rounded corners] (0,0)
  \foreach \i in {1,...,6} {
    -- ++(rnd,{rnd*.5}) node (cp) {}
};
    \node [circle,fill=black,inner sep=0pt,minimum size=1mm,opaque,label=right:$B$] at (cp) {};
    \node [circle,fill,inner sep=0pt,minimum size=1mm,opaque,label=left:$A$] at (0,0) {};


    \end{tikzpicture}
\caption{Траектория муравья}
\label{fig:p16}
\end{figure}

% \begin{figure}[H]
%     \centering

%     \includestandalone{PICS/mech/8ave_1}

% \caption{Траектория муравья}
% \label{fig:p16}
% \end{figure}


Муравей двигался по кривой траектории из пункта~$A$ в пункт~$B$ со случайно изменяющейся скоростью во времени, пройдя путь~$S_\text{общ}$ за время~$\Delta t_\text{общ}$. Если неинтересно, как тело двигалось в малые промежутки времени, то быстроту движения можно оценить по \textbf{средней скорости}:
\begin{equation}
    \label{8ave_e1}
    v_\text{ср}=\dfrac{S_\text{общ}}{\Delta t_\text{общ}}.
\end{equation}

Когда что-то известно о <<частях>> траектории (например, сколько участков выделяют в движении), то формулу~\eqref{8ave_e1} можно записать так:
\begin{equation*}
    v_\text{ср}=\dfrac{S_\text{1}+S_\text{2}+\ldots}{t_\text{1}+t_\text{2}+\ldots},
\end{equation*}
где $S_\text{1},S_\text{2}$~--- пути; $t_\text{1},t_\text{2}$~--- промежутки времени соответствующих участков.

Можно сказать, что средняя скорость~--- это такая скорость, которую нужно поддерживать телу, чтобы тот же путь был пройден за то же время, что и при действительном движении. 

Среднюю скорость, о которой говорилось выше, называют также \emph{путевой}. Этим подчеркивается, что во внимание принимается именно \emph{путь}~--- то есть длина какой-либо траектории. Теперь следует разобрать одну простую задачу.

\medskip

\noindent\textbf{Задача.} Санки, двигаясь по горе, прошли в течение первой секунды движения 2~м, второй секунды~--- 6~м, третьей секунды~--- 10~м и четвертой секунды~--- 14~м. Найдите среднюю скорость за первые две секунды, за последние две секунды и за все время.

\smallskip

\noindent\emph{Решение}. Движение <<разбито>> на 4~участка. Для первых двух можно найти:
\begin{equation*}
    v_\text{ср1,2}=\dfrac{S_\text{1}+S_\text{2}}{t_\text{1}+t_\text{2}}=\dfrac{2+6}{1+1}=4\ \text{м$/$с}.
\end{equation*}

Для вторых двух и всех четырех аналогично:
\begin{gather*}
    v_\text{ср3,4}=\dfrac{S_\text{3}+S_\text{4}}{t_\text{3}+t_\text{4}}=\dfrac{10+14}{1+1}=12\ \text{м$/$с},\\
    v_\text{ср1--4}=\dfrac{S_\text{1}+S_\text{2}+S_\text{3}+S_\text{4}}{t_\text{1}+t_\text{2}+t_\text{3}+t_\text{4}}=\dfrac{2+6+10+14}{1+1+1+1}=8\ \text{м$/$с}.
\end{gather*}

 
\end{document}